\chapter{Adjunctions}

\begin{overviewbox}
We have now built a substantial toolkit: categories, functors, natural
transformations, universal properties, limits, and colimits. The next
step is to study how categories themselves \emph{relate}.

An \term{adjunction} is the fundamental relationship between two
categories. It is the categorical formalization of a recurring informal
pattern: \emph{constructing} something on one side ``corresponds to''
\emph{checking} something on the other. The construction goes
\emph{left} (a functor $L$) and the checking goes \emph{right} (a functor
$R$), and the two are linked by a bijection between hom-sets:
\begin{equation*}
  \mathrm{Hom}_\mathcal{D}(L(A), B) \;\cong\; \mathrm{Hom}_\mathcal{C}(A, R(B)).
\end{equation*}
This bijection, when it holds naturally in both variables, is called an
\term{adjunction} and is written $L \dashv R$.

The reason adjunctions matter is that almost every important construction
in mathematics is part of an adjunction: free monoid/forgetful, free
group/forgetful, tensor/hom, $\exists$/pullback, abelianization/inclusion,
\ldots The list is essentially endless.

The deep fact of this chapter: left adjoints preserve colimits, right
adjoints preserve limits. This is what makes adjunctions the principal
tool for proving preservation theorems.
\end{overviewbox}


% ═══════════════════════════════════════════════════════════════════════════
\section{The Hom-Set Bijection}
\label{sec:hom-set-bijection}
% ═══════════════════════════════════════════════════════════════════════════

\subsection{Informally: Free and Forgetful}

Here is the prototypical pattern. Consider an object $A$ in some category
$\mathcal{C}$ of ``raw stuff'' (say a set). We want to ``freely build a
richer structure'' on $A$ --- a monoid, a group, a module --- landing in
some richer category $\mathcal{D}$. The free construction is $L(A)$, an
object of $\mathcal{D}$.

Conversely, given an object $B$ in $\mathcal{D}$, we can \emph{forget} its
extra structure, landing back in $\mathcal{C}$. The forgetful is $R(B)$,
an object of $\mathcal{C}$.

The magic: a structured map $L(A) \to B$ (i.e., a homomorphism in
$\mathcal{D}$) corresponds to an unstructured map $A \to R(B)$ (just a
function on raw stuff). Once you fix where each generator of $L(A)$ goes
in $B$, the rest of the homomorphism is forced. This correspondence is a
bijection $\mathrm{Hom}_\mathcal{D}(L(A), B) \cong \mathrm{Hom}_\mathcal{C}(A, R(B))$.

\subsection{Expanding: The Pattern in Concrete Settings}

\begin{example}{Free Monoid and Forgetful}
Take $L = \text{free monoid functor}, R = \text{forgetful from monoids to
sets}$. For a set $A$ and a monoid $B$:
\begin{itemize}
  \item A monoid homomorphism $L(A) \to B$ is determined by where it
        sends each $a \in A$ (the generators).
  \item A function $A \to R(B)$ is exactly a choice of an element of $B$
        for each $a \in A$.
\end{itemize}
These data are the same. The bijection is the formal statement of this
informal fact.
\end{example}

\begin{example}{Discrete and Forgetful for Topology}
Let $L$ send a set $X$ to ``$X$ with the discrete topology,'' and $R$
send a topological space to its underlying set. A continuous map
$L(X) \to Y$ is just any function $X \to R(Y)$, because every function
from a discrete space is continuous. So we get a bijection
\begin{equation*}
  \mathrm{Hom}_{\mathbf{Top}}(L(X), Y) \cong \mathrm{Hom}_{\mathbf{Set}}(X, R(Y)).
\end{equation*}
$L$ is left adjoint to $R$.
\end{example}

\begin{example}{Tensor and Hom}
For modules over a commutative ring $R$, the tensor product with a
fixed module $M$ is left adjoint to the hom-functor out of $M$:
\begin{equation*}
  \mathrm{Hom}(N \otimes M, P) \cong \mathrm{Hom}(N, \mathrm{Hom}(M, P)).
\end{equation*}
This bijection is so important it has a special name: \emph{tensor-hom
adjunction}. It generalizes the ``currying'' of functions in $\mathbf{Set}$:
$Z^{X \times Y} = (Z^Y)^X$.
\end{example}

\begin{dialog}
The pattern in each case: one functor \emph{builds something} and the
other \emph{forgets or evaluates}. The bijection says that the building
in one direction matches the using in the other.
\end{dialog}

\subsection{Formalizing: Adjoint Functors}

Let $L : \mathcal{C} \to \mathcal{D}$ and $R : \mathcal{D} \to \mathcal{C}$
be functors. We say $L$ is \term{left adjoint to} $R$ (equivalently $R$
is \term{right adjoint to} $L$), and we write $L \dashv R$, if there is a
\emph{natural} bijection
\begin{equation*}
  \varphi_{A,B} : \mathrm{Hom}_\mathcal{D}(L(A), B)
  \;\xrightarrow{\;\cong\;}\;
  \mathrm{Hom}_\mathcal{C}(A, R(B))
\end{equation*}
for every $A \in \mathcal{C}$ and $B \in \mathcal{D}$. The naturality
requires that $\varphi$ commute with precomposition by $L(f)$ in
$\mathcal{D}$ and with $f$ in $\mathcal{C}$, and similarly with $R(g)$.

If $\varphi$ exists, we say $L$ and $R$ form an \term{adjoint pair}, or
that $(L, R, \varphi)$ is an \term{adjunction}.

\begin{howtosay}
\begin{itemize}
  \item $L \dashv R$ --- ``$L$ is left adjoint to $R$,'' or
        ``the adjunction $L$-$R$.''
  \item The symbol $\dashv$ is read ``turnstile,'' written
        \texttt{\textbackslash dashv}. The reflected $\vdash$ goes the
        other direction.
  \item $\varphi$ is sometimes called the \emph{adjunction iso} or just
        \emph{the natural bijection}.
  \item Naturality is in \emph{both} variables: $A$ contravariantly
        (from $\mathcal{C}^{\mathrm{op}}$) and $B$ covariantly.
\end{itemize}
\end{howtosay}

\subsection{Reorganizing: Unit and Counit}

Plug $B = L(A)$ into the bijection: the right-hand side becomes
$\mathrm{Hom}_\mathcal{C}(A, R(L(A)))$. The identity $\mathrm{id}_{L(A)}$
on the left corresponds to a morphism on the right:
\begin{equation*}
  \eta_A : A \to R(L(A)).
\end{equation*}
This morphism is called the \term{unit of the adjunction at $A$}.

Dually, plug $A = R(B)$: the left-hand side becomes
$\mathrm{Hom}_\mathcal{D}(L(R(B)), B)$. The identity $\mathrm{id}_{R(B)}$
on the right corresponds to:
\begin{equation*}
  \varepsilon_B : L(R(B)) \to B,
\end{equation*}
the \term{counit at $B$}.

The collections $(\eta_A)$ and $(\varepsilon_B)$ assemble into natural
transformations
\begin{equation*}
  \eta : \mathrm{Id}_\mathcal{C} \Rightarrow R \circ L,
  \qquad
  \varepsilon : L \circ R \Rightarrow \mathrm{Id}_\mathcal{D}.
\end{equation*}

\subsection{Simplifying: The Triangle Identities}

\begin{align*}
  (R \varepsilon) \circ (\eta R) &= \mathrm{id}_R, \\
  (\varepsilon L) \circ (L \eta) &= \mathrm{id}_L.
\end{align*}
These are the \term{triangle identities} (or \emph{counit-unit laws}).
They encode the requirement that the bijection $\varphi$ is the natural
bijection of an adjunction, expressed entirely in terms of $L$, $R$,
$\eta$, and $\varepsilon$.

\subsection{Formally: Three Equivalent Definitions}

\formalnote{Adjunctions, three ways}
The following data are interchangeable:
\begin{enumerate}
  \item A natural bijection $\varphi$ as above.
  \item Natural transformations $\eta, \varepsilon$ satisfying the
        triangle identities.
  \item For every $A$, a universal arrow $\eta_A : A \to R(L(A))$ from
        $A$ to $R$ (or, dually, for every $B$, a universal arrow from
        $L$ to $B$).
\end{enumerate}
Each formulation is convenient in different proofs. The hom-set bijection
is best for computing examples; the unit/counit form is best for proofs
in 2-category theory; the universal-arrow form is best for the Adjoint
Functor Theorems (Chapter~10).


% ═══════════════════════════════════════════════════════════════════════════
\section{Adjoint Pairs in the Wild}
\label{sec:adjoints-in-wild}
% ═══════════════════════════════════════════════════════════════════════════

\subsection{Informally: A Tour}

The same hom-set bijection $\mathrm{Hom}(L A, B) \cong \mathrm{Hom}(A, R B)$
appears all over mathematics. Recognizing it is half the battle.

\subsection{Expanding: A Catalogue}

\begin{example}{Free Group --- Forgetful}
$L$ sends a set to the free group on it; $R$ sends a group to its
underlying set. The unit $\eta_A : A \to R L(A)$ is the inclusion of
generators.
\end{example}

\begin{example}{Abelianization --- Inclusion}
$L : \mathbf{Grp} \to \mathbf{Ab}$ sends a group $G$ to its abelianization
$G / [G, G]$. $R$ is the inclusion of abelian groups into all groups. The
unit is the quotient map $G \to G^{\mathrm{ab}}$.
\end{example}

\begin{example}{Discrete --- Forgetful for Posets}
$L$ sends a set $X$ to the discrete poset (no nontrivial order). $R$ sends
a poset to its underlying set. $L \dashv R$.
\end{example}

\begin{example}{Quantifiers as Adjoints}
For a function $f : X \to Y$, the pullback functor $f^* : \mathbf{Set}/Y
\to \mathbf{Set}/X$ has both a left adjoint $\exists_f$ and a right
adjoint $\forall_f$. The left adjoint is ``image'' and the right adjoint
is ``Frobenius's all-fiber condition.'' This makes existential and
universal quantification adjoints to substitution.
\end{example}

\begin{example}{Curry: Product --- Function-Space}
In $\mathbf{Set}$ (and any cartesian closed category):
\begin{equation*}
  \mathrm{Hom}(A \times B, C) \;\cong\; \mathrm{Hom}(A, C^B).
\end{equation*}
$(- \times B)$ is left adjoint to $(-)^B$.
\end{example}

\subsection{Formalizing: A Recipe for Identifying Adjunctions}

To spot an adjunction $L \dashv R$:
\begin{enumerate}
  \item Identify the two categories $\mathcal{C}$ and $\mathcal{D}$.
  \item Identify the candidate functors $L$ and $R$.
  \item Write out a typical hom-set $\mathrm{Hom}_\mathcal{D}(L A, B)$
        in concrete terms.
  \item Write out $\mathrm{Hom}_\mathcal{C}(A, R B)$ in concrete terms.
  \item Look for a natural bijection between the two descriptions.
\end{enumerate}
If you can carry out steps 3--5, you have your adjunction.

\subsection{Reorganizing: Adjoints Are Unique}

\textbf{Theorem (Uniqueness of adjoints).}\quad If $R$ has a left adjoint,
that left adjoint is unique up to natural isomorphism. Dually for right
adjoints.

The proof is the standard universal-property uniqueness argument from
\S5.4, applied to the unit (or counit).

\subsection{Simplifying: ``Adjoints'' Are Functorial Pairs}

We will routinely write things like ``the left adjoint to $R$ is $L$''
or ``$\exists_f$ is left adjoint to $f^*$.'' This is a statement about
a \emph{pair} of functors and a natural bijection, but the bijection is
canonically determined by either functor and need not be named.

\subsection{Formally: Composability of Adjunctions}

\formalnote{Adjunctions compose}
If $L \dashv R$ between $\mathcal{C}$ and $\mathcal{D}$, and $L' \dashv R'$
between $\mathcal{D}$ and $\mathcal{E}$, then $L' \circ L \dashv R \circ R'$
between $\mathcal{C}$ and $\mathcal{E}$. The bijection composes:
\begin{align*}
  \mathrm{Hom}_\mathcal{E}(L'L(A), B)
    &\cong \mathrm{Hom}_\mathcal{D}(L(A), R'(B)) \\
    &\cong \mathrm{Hom}_\mathcal{C}(A, R R'(B)).
\end{align*}


% ═══════════════════════════════════════════════════════════════════════════
\section{Adjunctions and Limits}
\label{sec:adjunctions-and-limits}
% ═══════════════════════════════════════════════════════════════════════════

\subsection{Informally: Right Adjoints Preserve Limits}

The most useful theorem about adjunctions is this:
\begin{quote}
\textit{If $R$ has a left adjoint, then $R$ preserves all limits that
exist in its source.}
\end{quote}
Dually, left adjoints preserve all colimits.

Why does this matter? Because limits are everywhere (Chapter~6) and
checking that a functor preserves them is, in general, painful. But once
we know our functor is a right adjoint, preservation is automatic.

\subsection{Expanding: Sketch of the Proof}

The proof is short. Suppose $L \dashv R$, $D : \mathcal{J} \to \mathcal{D}$
is a diagram, and $\lim D$ exists. We want to show $R(\lim D) = \lim (R
\circ D)$.

For any object $A$:
\begin{align*}
  \mathrm{Hom}(A, R(\lim D))
    &\cong \mathrm{Hom}(L(A), \lim D) && \text{(adjunction)} \\
    &\cong \lim_j \mathrm{Hom}(L(A), D(j)) && \text{(limit/hom)} \\
    &\cong \lim_j \mathrm{Hom}(A, R(D(j))) && \text{(adjunction, pointwise)} \\
    &\cong \mathrm{Hom}(A, \lim_j R(D(j))). && \text{(limit/hom)}
\end{align*}
This natural isomorphism for every $A$ identifies $R(\lim D)$ with $\lim
(R \circ D)$ up to unique iso. Hence $R$ preserves the limit.

\subsection{Formalizing: The Statement, Precisely}

\textbf{Theorem (RAPL: Right Adjoints Preserve Limits).}\quad Let
$L \dashv R$ with $L : \mathcal{C} \to \mathcal{D}$ and $R : \mathcal{D}
\to \mathcal{C}$. For any diagram $D : \mathcal{J} \to \mathcal{D}$ whose
limit exists, the canonical morphism
\begin{equation*}
  R(\lim D) \;\xrightarrow{\;\cong\;}\; \lim (R \circ D)
\end{equation*}
is an isomorphism.

Dually: \textbf{LAPC (Left Adjoints Preserve Colimits).}

\subsection{Reorganizing: A Test for ``No Left Adjoint''}

The contrapositive is useful: if $R$ \emph{fails} to preserve some limit,
then $R$ has no left adjoint. For example:
\begin{itemize}
  \item The functor $\pi_0 : \mathbf{Top} \to \mathbf{Set}$ (path
        components) fails to preserve products (cf. the line $\times$
        line example), so it has no left adjoint (in the unenriched
        sense).
  \item The free-group functor $L : \mathbf{Set} \to \mathbf{Grp}$ does
        not preserve products, so it has no right adjoint either. (It is
        itself a left adjoint, of course --- of the forgetful functor.)
\end{itemize}

\subsection{Simplifying: The Adjunction Slogan}

\begin{equation*}
  \boxed{\;\;
    L \dashv R
    \;\Rightarrow\;
    R \text{ preserves limits},\;\;
    L \text{ preserves colimits}.
  \;\;}
\end{equation*}

\subsection{Formally: Hom-Functor Preservation}

\formalnote{Hom-functors and limits}
For any object $A$ of $\mathcal{C}$, the hom-functor
$\mathrm{Hom}(A, -)$ preserves limits. This is sometimes called
``co-Yoneda preservation'' and it is the limit/hom step used in the
proof of RAPL. Dually, $\mathrm{Hom}(-, A)$ (contravariant) sends
colimits to limits.


% ═══════════════════════════════════════════════════════════════════════════
\section{Equivalences as Adjunctions}
\label{sec:equivalences}
% ═══════════════════════════════════════════════════════════════════════════

\subsection{Informally: When the Two Sides Are Almost the Same}

A special and important case: when both the unit $\eta$ and the counit
$\varepsilon$ are natural isomorphisms. Then $L$ and $R$ are inverses up
to natural isomorphism, and the two categories $\mathcal{C}$ and
$\mathcal{D}$ are essentially the same. We call this an
\term{equivalence} of categories.

\subsection{Expanding: Equivalence vs. Isomorphism}

An \emph{isomorphism} of categories would require $R \circ L = \mathrm{Id}_\mathcal{C}$
and $L \circ R = \mathrm{Id}_\mathcal{D}$ \emph{on the nose}. This is too
strict --- and usually irrelevant. What we really want is that the two
sides are the same \emph{up to natural isomorphism}, which is exactly
what equivalence demands.

\begin{example}{Finite Sets and Their Cardinality Skeletons}
Let $\mathcal{C}$ be $\mathbf{FinSet}$ (all finite sets and functions),
and $\mathcal{D}$ be the skeleton $\mathbf{FinSet}_\mathrm{sk}$ with
objects $\underline{n} = \{1, 2, \ldots, n\}$ for $n \geq 0$. The
inclusion $\mathcal{D} \hookrightarrow \mathcal{C}$ is one half of an
equivalence; the other functor sends a finite set $X$ to
$\underline{|X|}$. Every finite set is isomorphic to some $\underline{n}$,
so the unit and counit are natural isomorphisms.
\end{example}

\subsection{Formalizing: Equivalence of Categories}

An \term{equivalence of categories} between $\mathcal{C}$ and $\mathcal{D}$
is a pair of functors $L : \mathcal{C} \to \mathcal{D}$ and
$R : \mathcal{D} \to \mathcal{C}$ together with natural isomorphisms
\begin{equation*}
  R \circ L \;\cong\; \mathrm{Id}_\mathcal{C},
  \qquad
  L \circ R \;\cong\; \mathrm{Id}_\mathcal{D}.
\end{equation*}
We write $\mathcal{C} \simeq \mathcal{D}$.

Equivalently: $L \dashv R$, both unit and counit are natural isomorphisms.

\textbf{Theorem (Recognition of Equivalences).}\quad A functor
$F : \mathcal{C} \to \mathcal{D}$ is one half of an equivalence iff $F$
is fully faithful and \term{essentially surjective} (every object of
$\mathcal{D}$ is isomorphic to some $F(C)$).

\subsection{Reorganizing: The Strict-Equivalence Slogan}

Isomorphism of categories is the strict version; equivalence is the
\emph{right} version. Whenever we say ``two categories are the same'' we
almost always mean equivalence.

\subsection{Simplifying: The Equivalence Schema}

\begin{equation*}
  F \text{ is an equivalence}
  \;\iff\;
  F \text{ is fully faithful and essentially surjective}.
\end{equation*}

\subsection{Formally: Working with Skeletons}

\formalnote{Skeletons and equivalence}
Every category $\mathcal{C}$ is equivalent to its \term{skeleton}: a
subcategory containing one object per isomorphism class. Working with a
skeleton is often convenient, because it eliminates redundant objects
without changing the categorical content.


% ═══════════════════════════════════════════════════════════════════════════
\section*{Exercises}
\addcontentsline{toc}{section}{Exercises}
% ═══════════════════════════════════════════════════════════════════════════

\begin{exercisesbox}
\begin{qa}
\qaitem{What is an adjunction, in one sentence?}{A natural bijection between $\mathrm{Hom}_\mathcal{D}(L A, B)$ and $\mathrm{Hom}_\mathcal{C}(A, R B)$.}
\qaitem{Which is the left adjoint: $L$ or $R$?}{$L$. It appears on the left of the hom inside $\mathcal{D}$.}
\qaitem{What does $L \dashv R$ mean?}{$L$ is left adjoint to $R$.}
\qaitem{In ``free monoid --- forgetful,'' which is left?}{Free monoid is left; forgetful is right.}
\qaitem{What is the unit of an adjunction?}{A natural transformation $\eta : \mathrm{Id}_\mathcal{C} \Rightarrow R L$.}
\qaitem{What is the counit?}{A natural transformation $\varepsilon : L R \Rightarrow \mathrm{Id}_\mathcal{D}$.}
\qaitem{How does $\eta_A$ relate to $\mathrm{id}_{L(A)}$?}{Under the bijection, $\eta_A$ corresponds to $\mathrm{id}_{L(A)}$.}
\qaitem{How does $\varepsilon_B$ relate to $\mathrm{id}_{R(B)}$?}{Under the inverse bijection, $\varepsilon_B$ corresponds to $\mathrm{id}_{R(B)}$.}
\qaitem{What are the triangle identities?}{$(R \varepsilon)(\eta R) = \mathrm{id}_R$ and $(\varepsilon L)(L \eta) = \mathrm{id}_L$.}
\qaitem{Why do we care about the triangle identities?}{They are the data-only form of the adjunction (no bijection needed).}
\qaitem{Are adjoints unique?}{Yes, up to natural isomorphism.}
\qaitem{What is the slogan for limits and adjoints?}{Right adjoints preserve limits; left adjoints preserve colimits. (RAPL/LAPC.)}
\qaitem{If $R$ fails to preserve some limit, what can you conclude?}{$R$ has no left adjoint.}
\qaitem{Why does the hom-set bijection imply RAPL?}{Apply $\mathrm{Hom}(A, -)$, which preserves limits; chase the bijection.}
\qaitem{What is the prototype adjunction in $\mathbf{Set}$?}{Currying: $(- \times B) \dashv (-)^B$.}
\qaitem{What is the tensor-hom adjunction in modules?}{$(- \otimes M) \dashv \mathrm{Hom}(M, -)$.}
\qaitem{Quantifiers as adjoints to what?}{Adjoints to the pullback functor $f^*$: $\exists_f \dashv f^* \dashv \forall_f$.}
\qaitem{What is an equivalence of categories?}{A pair $L \dashv R$ where both unit and counit are natural isomorphisms.}
\qaitem{Why prefer equivalence to isomorphism of categories?}{Isomorphism is too strict; almost all ``sameness'' in CT is up to natural iso.}
\qaitem{What is essential surjectivity?}{Every object in $\mathcal{D}$ is isomorphic to some $F(C)$.}
\qaitem{When is $F$ an equivalence?}{When $F$ is fully faithful and essentially surjective.}
\qaitem{What is a skeleton of a category?}{A subcategory with one object per isomorphism class.}
\qaitem{Is every category equivalent to its skeleton?}{Yes.}
\qaitem{Do adjunctions compose?}{Yes: $L \dashv R$ and $L' \dashv R'$ give $L' L \dashv R R'$.}
\qaitem{Does the identity functor have an adjoint?}{Yes; it is adjoint to itself.}
\qaitem{Can a functor have more than one left adjoint?}{Only up to natural isomorphism; there is essentially one.}
\qaitem{Why is the unit ``the obvious'' inclusion in many examples?}{Because the unit corresponds to $\mathrm{id}_{L(A)}$, the most basic morphism on the freely built object.}
\qaitem{What kind of functor is $\pi_0 : \mathbf{Top} \to \mathbf{Set}$?}{Not a right adjoint in the unenriched sense (it fails to preserve some products).}
\qaitem{What is the practical use of recognizing an adjunction?}{You instantly get a wealth of preservation theorems and natural bijections.}
\qaitem{What's the deeper meaning of $L \dashv R$?}{The construction $L$ \emph{represents} the same data as the checking $R$, viewed dually.}
\end{qa}
\end{exercisesbox}


% ═══════════════════════════════════════════════════════════════════════════
\section*{Chapter Summary}
\addcontentsline{toc}{section}{Chapter Summary}
% ═══════════════════════════════════════════════════════════════════════════

\begin{summarybox}
An \textbf{adjunction} $L \dashv R$ is a natural bijection
$\mathrm{Hom}(L A, B) \cong \mathrm{Hom}(A, R B)$. Equivalently, it is a
pair of natural transformations $\eta : \mathrm{Id} \Rightarrow R L$
(\textbf{unit}) and $\varepsilon : L R \Rightarrow \mathrm{Id}$
(\textbf{counit}) satisfying the \textbf{triangle identities}.

Adjoints are unique up to natural isomorphism. They compose. They are
ubiquitous: free/forgetful pairs, tensor-hom, quantifiers, abelianization,
discrete topology, etc.

The major theorem: \textbf{Right Adjoints Preserve Limits} (RAPL) and
\textbf{Left Adjoints Preserve Colimits} (LAPC).

An \textbf{equivalence of categories} is an adjunction with invertible
unit and counit. It is recognized by being \textbf{fully faithful and
essentially surjective}.
\end{summarybox}


% ═══════════════════════════════════════════════════════════════════════════
\section*{Formal Restatement}
\addcontentsline{toc}{section}{Formal Restatement}
% ═══════════════════════════════════════════════════════════════════════════

\begin{formalrestatement}
\textbf{Adjunction.}\quad $L : \mathcal{C} \to \mathcal{D}$ is left
adjoint to $R : \mathcal{D} \to \mathcal{C}$ ($L \dashv R$) iff there is a
natural bijection
\begin{equation*}
  \varphi_{A,B} : \mathrm{Hom}_\mathcal{D}(L A, B) \cong \mathrm{Hom}_\mathcal{C}(A, R B).
\end{equation*}

\medskip
\textbf{Unit and counit.}\quad
$\eta : \mathrm{Id}_\mathcal{C} \Rightarrow R L$,
$\varepsilon : L R \Rightarrow \mathrm{Id}_\mathcal{D}$, satisfying:
\begin{equation*}
  (R \varepsilon)(\eta R) = \mathrm{id}_R,
  \qquad
  (\varepsilon L)(L \eta) = \mathrm{id}_L.
\end{equation*}

\medskip
\textbf{Limit preservation.}\quad
$L \dashv R \;\Rightarrow\; R \text{ preserves limits},\; L \text{ preserves colimits}$.

\medskip
\textbf{Equivalence.}\quad
$\mathcal{C} \simeq \mathcal{D}$ iff there exist $L \dashv R$ with
$\eta, \varepsilon$ both natural isomorphisms; equivalently, there exists
a fully faithful, essentially surjective functor.

\medskip
\noindent\textit{Chapter~8 develops monads, which arise from adjunctions
and capture algebraic structure relative to a base category.}
\end{formalrestatement}
